% commands.tex
% Definiciones de comandos personalizados compartidos por todo el documento.
% Se carga una sola vez desde el preámbulo de main.tex.
%
% Se usa \providecommand (en lugar de \newcommand) para que la definición no
% entre en conflicto si algún paquete ya define el mismo nombre.

% Marcador de decisión pendiente. Visible en el PDF compilado para no olvidar
% completarlo antes de la entrega.
\providecommand{\todo}[1]{\textbf{[TODO: #1]}}

% Stub de figura: incluye la imagen si el archivo existe y, mientras no exista,
% dibuja un recuadro de marcador de posición para que el documento siga
% compilando. Requiere graphicx y float (cargados en main.tex).
\providecommand{\figstub}[2][\linewidth]{%
  \IfFileExists{#2}{%
    \includegraphics[width=#1]{#2}%
  }{%
    \fbox{\parbox[c][4.5cm][c]{0.82\linewidth}{%
      \centering\itshape [Figura pendiente: \texttt{\detokenize{#2}}]}}%
  }%
}
